% Integrated core working draft; final independent-source research and release QA pending.
% Evidence scope: original packet controls and actual full-call service; no encoder-cache result.
\subsection{Native feedback and conditional full calls}
\label{subsec:native_full_calls}
StreamCorrect's original recognizer, learned correction, and feedback update
provide a concrete instance of the contract in Section~\ref{sec:contract}
\cite{vu2026streamcorrect}. Let $P$ be a complete closed native state and
$X$ the next accepted audio input. We expose the recognition boundary as
\begin{equation}
  A(P,X)=(B,M),
\end{equation}
where $B$ contains the ordered beam candidates and $M$ contains the live
engine state and inputs required by the remaining correction and update.
The protected original $C/U$ selects a correction and produces the next
closed state. Its selection, candidate order, prompts, returned token IDs,
and update rules remain intact.

Before that selection finishes, the original update tail supplied with a
typed first-beam identity prediction produces a provisional closed state
$\widehat P^+$. This projection has explicit provenance: recognition has
executed, but learned correction has not completed for this projected node.
It supplies a conditional input for the next complete call. Both GPUs host
the original full recognizer--corrector pipeline, so a call dispatched on
this input executes the protected original $C/U$ as well. Completed
conditional full results and typed identity projections remain distinct
objects. Neither is installed merely because its predicted text looks
plausible.

When the selected parent arrives, the Owner compares the call's entire
frozen start state, actual frontier, execution contract, and session epoch
with the installed parent. An exact match confirms the already owned call:
an in-flight call continues, or an already completed result becomes eligible
for ordered installation. The job's birth root, identifier, epoch, inputs,
and execution history are unchanged. The call's own identity prediction need
not equal its eventual corrected output; confirming its input establishes
eligibility, while original correction determines its result.

An earlier ancestor mismatch removes support for the tracked provisional
chain, but does not prove that a later complete input cannot converge.
Cancellation before correction therefore requires an actual selected state
at the call's same start frontier that disagrees with its complete input,
coverage of its target without remaining use, or a changed epoch. A still
unresolved future parent supplies none of these proofs. Loss of current
view support can therefore require withdrawal without yet authorizing
cancellation of that unresolved call.

This direct native witness instantiates the continuation relation of
Section~\ref{sec:continuation}. It differs from the core Qwen computation
guard, which conservatively requires an explicit matching boundary node and
a uniquely tracked successor. These are different sufficient implementations
of eligibility. The whole-view Consumer likewise keeps certified and
provisional provenance separate and applies complete supported hypotheses.
The native service comparison uses two full-pipeline workers in both arms
and differs in conditional dispatch and retention. Ready-only ASR handoff
is a separate developmental control. These are author-built implementations
of the feedback principles; the measured benefit concerns confirmed full calls.
